\documentclass[10pt]{extarticle}
\usepackage[a4paper,margin=1.5cm]{geometry}
\usepackage[T1]{fontenc}
\usepackage{amsmath,amssymb,mathtools}
\usepackage{graphicx}
\usepackage[export]{adjustbox}
\usepackage{multicol}
\usepackage{enumitem}
\usepackage{booktabs,array,tabularx}
\setlength{\columnsep}{0.7cm}
\setlength{\parindent}{0pt}
\setlength{\parskip}{3pt}
\pagestyle{empty}
\makeatletter
\newcommand{\qfigdef}[4]{\@namedef{qf@file@#1}{#2}\@namedef{qf@w@#1}{#3}\@namedef{qf@wide@#1}{#4}}
\newcommand{\qf@check}[1]{\@ifundefined{qf@file@#1}{\PackageError{score2}{Unknown figure id #1}{}}{}}
\newcommand{\qf@img}[4][]{\includegraphics[width=#2\dimexpr\@nameuse{qf@w@#4}pt\relax,max width=#3,max height=0.45\textheight,#1]{\@nameuse{qf@file@#4}}}
\newif\ifqf@full
\newif\ifqf@boxed
\newsavebox{\qf@box}
\newcommand{\qf@open}{\global\qf@fulltrue\par\global\qf@boxedtrue
  \setbox\qf@box\vbox\bgroup\hsize\textwidth\linewidth\textwidth}
\newcommand{\qf@emit}[1]{\par\egroup\global\qf@boxedfalse
  \ifdim\dimexpr\ht\qf@box+\dp\qf@box\relax<#1\textheight\noindent\usebox{\qf@box}\par\else\unvbox\qf@box\fi}
\newcommand{\qfullsplit}{\ifqf@boxed\qf@emit{0.85}\fi}
\newcommand{\qf@close}{\ifqf@boxed\qf@emit{0.45}\fi\global\qf@fullfalse}
\newcommand{\qfullbegin}{\end{multicols}\qf@open}
\newcommand{\qfullend}{\qf@close\begin{multicols}{2}}
\newcommand{\qfullstart}{\qf@open}
\newcommand{\qfulldone}{\qf@close}
\newcommand{\qfig}[2][1]{\qf@check{#2}%
  \ifqf@full\qf@blk{#1}{#2}\else
  \ifnum\@nameuse{qf@wide@#2}=1 \expandafter\@firstoftwo\else\expandafter\@secondoftwo\fi
  {\end{multicols}\begin{center}\qf@img{#1}{\textwidth}{#2}\end{center}\begin{multicols}{2}}%
  {\qf@blk{#1}{#2}}\fi}
\newcommand{\qf@blk}[2]{\par\smallskip{\centering\qf@img{#1}{\linewidth}{#2}\par}\smallskip}
\newcommand{\qfigi}[2][1]{\qf@check{#2}\qf@img[valign=c]{#1}{\linewidth}{#2}}
\makeatother
\newcommand{\papersec}[1]{\par\bigskip{\centering\Large\bfseries PAPER-#1\par}\medskip}
\newcommand{\partsec}[2]{\par\medskip\noindent{\large\bfseries PART-#1: #2}\par\smallskip}
\newcommand{\sect}[1]{\par\smallskip\noindent\textbf{#1}\par}
\newcommand{\qn}[1]{\par\smallskip\noindent\textbf{#1.}~}
\newcommand{\paraq}[2]{\par\smallskip\noindent\textbf{Paragraph for Questions #1 and #2}\par}
\newenvironment{qoptions}{\begin{itemize}[leftmargin=1.8em,itemsep=1pt,topsep=2pt,label={}]}{\end{itemize}}
\newcommand{\opt}[2]{\item[(#1)] #2}
\newcolumntype{L}[1]{>{\raggedright\arraybackslash}p{#1}}
\newcommand{\fwstart}{\end{multicols}\par\noindent\begin{minipage}{\textwidth}\small\setlength{\tabcolsep}{4pt}\renewcommand{\arraystretch}{1.5}}
\newcommand{\fwend}{\end{minipage}\par\medskip\begin{multicols}{2}}
% figure registry (generated by score2_tools.py finalize)
\qfigdef{p1_phy_s1_q01_1}{SRG (Enthusiast) - Test 14 (21 Nov 2021)-img/p1_phy_s1_q01_1.png}{95.1}{0}
\qfigdef{p1_phy_s1_q02_1}{SRG (Enthusiast) - Test 14 (21 Nov 2021)-img/p1_phy_s1_q02_1.png}{107.5}{0}
\qfigdef{p1_phy_s1_q03_1}{SRG (Enthusiast) - Test 14 (21 Nov 2021)-img/p1_phy_s1_q03_1.png}{174.7}{0}
\qfigdef{p1_phy_s1_q04_1}{SRG (Enthusiast) - Test 14 (21 Nov 2021)-img/p1_phy_s1_q04_1.png}{196.8}{0}
\qfigdef{p1_phy_s2_q09_1}{SRG (Enthusiast) - Test 14 (21 Nov 2021)-img/p1_phy_s2_q09_1.png}{119.4}{0}
\qfigdef{p1_phy_s2_q09_2}{SRG (Enthusiast) - Test 14 (21 Nov 2021)-img/p1_phy_s2_q09_2.png}{105.2}{0}
\qfigdef{p1_phy_s2_q09_3}{SRG (Enthusiast) - Test 14 (21 Nov 2021)-img/p1_phy_s2_q09_3.png}{150.5}{0}
\qfigdef{p1_phy_s2_q09_4}{SRG (Enthusiast) - Test 14 (21 Nov 2021)-img/p1_phy_s2_q09_4.png}{158.4}{0}
\qfigdef{p1_phy_s2_q10_1}{SRG (Enthusiast) - Test 14 (21 Nov 2021)-img/p1_phy_s2_q10_1.png}{115.2}{0}
\qfigdef{p1_phy_s3_q12_1}{SRG (Enthusiast) - Test 14 (21 Nov 2021)-img/p1_phy_s3_q12_1.png}{121.3}{0}
\qfigdef{p1_phy_s3_q13_1}{SRG (Enthusiast) - Test 14 (21 Nov 2021)-img/p1_phy_s3_q13_1.png}{99.5}{0}
\qfigdef{p1_phy_s3_q14_1}{SRG (Enthusiast) - Test 14 (21 Nov 2021)-img/p1_phy_s3_q14_1.png}{92.1}{0}
\qfigdef{p1_phy_s3_q15_2}{SRG (Enthusiast) - Test 14 (21 Nov 2021)-img/p1_phy_s3_q15_2.png}{71.8}{0}
\qfigdef{p1_phy_s3_q16_1}{SRG (Enthusiast) - Test 14 (21 Nov 2021)-img/p1_phy_s3_q16_1.png}{271.0}{1}
\qfigdef{p1_phy_s3_q17_1}{SRG (Enthusiast) - Test 14 (21 Nov 2021)-img/p1_phy_s3_q17_1.png}{141.6}{0}
\raggedcolumns\makeatletter\@beginparpenalty=10000\makeatother

\begin{document}
\begin{center}{\Large\bfseries SRG (Enthusiast) --- Test 14 (21 Nov 2021)}\end{center}
\begin{multicols}{2}
\partsec{1}{PHYSICS}
\sect{SECTION-I}
\qn{1} Figure shows non-uniform magnetic field developed by magnetic rod. The magnetic field is described in terms of $B_r=0.3r,\,B_z=(1-2z)$ where $r$, $z$ are radial and vertical position coordinates respectively. A resistance less ring of mass $m$ radius $a$, inductance $L$ is shown in figure. Initially centre of ring is present at $r=0$, $z=0$ and Current flowing in it is zero. ($g$ is acceleration due to gravity)
\qfig{p1_phy_s1_q01_1}
\begin{qoptions}
\opt{A}{Current through ring is $\dfrac{2z\pi a^2}{L}$}
\opt{B}{Motion of ring is periodic but not S.H.M}
\opt{C}{Maximum $z$ coordinate of centre of ring during its motion is $\dfrac{mgL}{0.6\pi^2a^4}$}
\opt{D}{Amplitude of the motion is $\dfrac{mgL}{1.2\pi^2a^4}$}
\end{qoptions}
\qn{2} A plane electromagnetic sinusoidal wave propagating in the $+x$ direction. At $x=a$ electric field is $E_0(-\hat j)$ and at $x=a+dx$ it is $(E_0+dE)(-\hat j)$ then ($B$ -- magnetic field)
\qfig{p1_phy_s1_q02_1}
\begin{qoptions}
\opt{A}{Direction of $B$ in shaded region is $(-\hat k)$ and rate of change of $B$ with respective time is $+$ve}
\opt{B}{Direction of $B$ in shaded region is $(+\hat k)$ and rate of change of $B$ with respective time is $-$ve}
\opt{C}{Direction of $B$ in shaded region is $(+\hat k)$ and rate of change of $B$ with respective time is $+$ve}
\opt{D}{Direction of $B$ in shaded region is $(-\hat k)$ and rate of change of $B$ with respective time is $-$ve}
\end{qoptions}
\qn{3} In shown circuit rms voltage across inductor and resistor is 100 V
\qfig{p1_phy_s1_q03_1}
\begin{qoptions}
\opt{A}{$V_{rms}$ can be any value between 0 to $100\sqrt2$ V}
\opt{B}{If $V_{rms}=100$ V and $I_R=\dfrac{5\sqrt3}{2}$ A then $I=5$ A}
\opt{C}{If $V_{rms}=100$ V and $I_R=\dfrac{5\sqrt3}{2}$ A then capacitive reactance is $40\,\Omega$}
\opt{D}{$V_{rms}$ can be any value between 0 to 200 V}
\end{qoptions}
\qn{4} A circuit shown in figure contains ideal inductor and battery. Initially current through inductor is zero and the switch is closed at $t=0$. When potential difference across inductor is 50\% of maximum potential difference across it then select correct alternative/s:
\qfig{p1_phy_s1_q04_1}
\begin{qoptions}
\opt{A}{Current through battery is $7E/12R$}
\opt{B}{Current through resistance $R_2$ is $5E/12R$}
\opt{C}{Current through inductor is $E/6R$}
\opt{D}{Potential difference across $R_1$ is $E/2$}
\end{qoptions}
\qn{5} The two slits of double slit experiment were moved symmetrically apart with small relative velocity $v$ and the distance between screen and mid-point of slits is fixed and equal to $D$. Consider a point $P$ on the screen at a distance $x$ from central maxima then $(x\lll D)$:
\begin{qoptions}
\opt{A}{Rate of change of number of fringes between central maxima and point $P$ changes with respect to time is $\dfrac{xv}{\lambda D}$.}
\opt{B}{Number of fringes contained between central maxima and point $P$ increases with time.}
\opt{C}{Fringe width decreases as time passes.}
\opt{D}{Maxima's are moving closer to each other.}
\end{qoptions}
\qn{6} In the nuclei $X$ decay series begins with $^{232}X$. (Here 232 is mass number). $^{228}X$ alpha - decays into $^{224}Y$ with a half of 10000 seconds, $^{224}Y$ alpha-decays into $^{220}Z$ with a half life of 100 seconds, which in turn decays into $^{216}W$ with a half life of 50 s. If the nuclides are in radioactive equilibrium in a mineral sample that contains 1.0 g of $^{220}Z$. At this instant
\begin{qoptions}
\opt{A}{Mass of $^{232}$Th nearly is 310 g}
\opt{B}{Mass of $^{220}$Rn nearly is 2 g}
\opt{C}{Mass of $^{232}$Th nearly is 200 g}
\opt{D}{Mass of $^{220}$Rn nearly is 200 g}
\end{qoptions}
\qn{7} The activity of a sample of radioactive material is $R_1$ at time $t_1$ and $R_2$ at time $t_2$ $(t_2>t_1)$. If mean life of radioactive sample is $T$, then
\begin{qoptions}
\opt{A}{$R_1t_1=R_2t_2$}
\opt{B}{$\dfrac{R_1-R_2}{t_2-t_1}=\text{constant}$}
\opt{C}{$R_2=R_1\exp\!\left(\dfrac{t_2-t_1}{T}\right)$}
\opt{D}{$R_2=R_1\exp\!\left(\dfrac{t_1}{Tt_2}\right)$}
\end{qoptions}
\qn{8} A moving H-atom makes a head on collision with a stationary Li$^{++}$ ion. Just before collision, both H and Li$^{++}$ are in their first excited state. Immediately after collision H was found to be in ground state and Li$^{++}$ was in second excited state. No photon is emitted in the process. Assume that mass of Li$^{++}$ is nearly 7 times the mass of H.
\begin{qoptions}
\opt{A}{Initial kinetic energy of H-atom can be 7 eV}
\opt{B}{Initial kinetic energy of H-atom can be 10 eV}
\opt{C}{Initial kinetic energy of H-atom can be 6 eV}
\opt{D}{Initial kinetic energy of H-atom can be 8.7 eV}
\end{qoptions}
\fwstart
\sect{SECTION-IV}
\qn{9} Match the following\par\nobreak
\begin{center}\begin{tabular}{|L{0.6cm}|L{6.8cm}|L{0.8cm}|L{4.8cm}|}\hline
\multicolumn{2}{|c|}{\textbf{Column I}} & \multicolumn{2}{c|}{\textbf{Column II}}\\\hline
(A) & An ideal mono atomic gas undergoes the following cyclic process.

\qfigi[0.8]{p1_phy_s2_q09_1} & (P) & Work done in process $A$ to $B$ is 0\\\hline
(B) & An ideal polyatomic non-linear gas undergoes the cyclic process shown in (Take : $2^{4/3}=2.5$)

\qfigi[0.8]{p1_phy_s2_q09_2} & (Q) & Change in internal energy from $B$ to $C$ is 0\\\hline
(C) & An ideal monoatomic gas undergoes the following process.

\qfigi[0.8]{p1_phy_s2_q09_3} & (R) & Efficiency of the cycle is less than 0.5.\\\hline
(D) & An ideal mono atomic gas undergoes the following cyclic process. The pressure at point $A$ is $P_0$

\qfigi[0.8]{p1_phy_s2_q09_4} & (S) & Work done in process $A$ to $B$ is $(3/4)P_0V_0$\\\hline
 & & (T) & Heat supplied to the gas in process $A$ to $B$ is positive\\\hline
\end{tabular}\end{center}
\fwend
\qn{10} A monochromatic parallel beam of light of wavelength $\lambda$ is incident normally on the plane containing slits $S_1$ and $S_2$. The slits are of unequal width such that intensity only due to one slit on screen is four times that only due to the other slit. The screen is placed perpendicular to $x$-axis as shown. The distance between slits is $d$ and that between screen and slit is $D$. Match the statements in Column - I with results in Column - II.
\qfig{p1_phy_s2_q10_1}
\fwstart
\begin{center}\begin{tabular}{|L{0.6cm}|L{7.4cm}|L{0.8cm}|L{2.4cm}|}\hline
\multicolumn{2}{|c|}{\textbf{Column - I}} & \multicolumn{2}{c|}{\textbf{Column - II}}\\\hline
(A) & The distance between two points on screen having equal intensities, such that intensity at those points is $\dfrac19$th of maximum intensity & (P) & $\dfrac{D\lambda}{3d}$\\\hline
(B) & The distance between two points on screen having equal intensities, such that intensity at those points is $\dfrac39$th of maximum intensity & (Q) & $\dfrac{D\lambda}{d}$\\\hline
(C) & The distance between two points on screen having equal intensities, such that intensity at those points is $\dfrac59$th of maximum intensity & (R) & $\dfrac{2D\lambda}{d}$\\\hline
(D) & The distance betwen two points on screen having equal intensities, such that intensity at those points is $\dfrac79$th of maximum intensity. & (S) & $\dfrac{3D\lambda}{d}$\\\hline
 & & (T) & $\dfrac{2D\lambda}{3d}$\\\hline
\end{tabular}\end{center}
\fwend
\sect{SECTION-II}
\qn{11} A box containing an inductor of inductance $L$, a capacitor of capacitance $C$ and a resistor of resistance $R$. When 250 V d.c. is applied to the terminals of the box, a current of 1.0 A plus an oscillating current flows in the current. When an a.c. source of 250 V (r.m.s) at 2250 rad/sec is connected, an rms current of 1.25 A is produced. This current rises with frequency and becomes maximum at 4500 rad/sec. If $2R\sqrt{CL}$ is $x$ SI unit then $9x$ is \underline{\hspace{1.5cm}}
\qn{12} In $L$-$C$ circuit shown in the figure, the value of current at $t=0$ is $i_0$ and charge on the capacitor at the same moment is $i_0\sqrt{3LC}$. The value of the maximum current through the inductor is $\dfrac n2 i_0$, then $n$ is
\qfig{p1_phy_s3_q12_1}
\qn{13} A series $R$-$C$ circuit is connected to an AC supply as shown in the figure. If power factor of the circuit is $\dfrac1{\sqrt5}$, then the value of the capacitance $C$ (in $\mu$F) is $\dfrac{1000}{n}$ then $n$ is
\qfig{p1_phy_s3_q13_1}
\qn{14} Two long solenoids are arranged coaxially, there radii are $R$ and $2R$, there coils have same number of turns per unit length. Current starts from 0 and increase at the rate $3k$ and $2k$ (in both coils in same sense) inner and outer coils respectively. As a result of increasing currents charged particle initially at rest between the solenoids it starts moving along a circular trajectory of radius $\sqrt{a/b}$ then find $(a+b)$.
\qfig{p1_phy_s3_q14_1}
\qn{15} A conducting disc of mass `$m$' and volume `$V$' is suspended with the help of a light spring of force constant `$k$' from a fixed support. Thickness of the disc is much smaller than its radius. A uniform magnetic field of induction `$B$' parallel to the plane of the disc is established. The disc is pulled down by very small distance $b$ from equilibrium position and then released. $\left(B^2V\varepsilon_0=3M\right)$. Then maximum velocity of the disc is $\dfrac bn\sqrt{\dfrac km}$. Find $n$.
\qfig{p1_phy_s3_q15_2}
\qfullbegin
\qn{16} The inductors of two $LR$ circuits are placed next to each other, as shown in the figure. The values of self-inductance of the inductors, resistances, mutual-inductance and applied voltages are specified in the given circuit. After both the switches are closed simultaneously, the total work done by the batteries against the induced EMF in the inductors by the time the currents reach their steady values is $x$ mJ then $\dfrac x5$ is \underline{\hspace{1.5cm}}
\qfig{p1_phy_s3_q16_1}
\qfullend
\qn{17} A cyclic process (cycle) 1-2-3-4-1 consisting of two isobars 2-3 and 4-1, isochor 1-2, and a certain process 3-4 represented by a straight line on the $P$-$V$ diagram (see Fig.) involves $n$ moles of an ideal gas. The gas temperatures in states 1, 2 and 3 are $T_1$, $T_2$ and $T_3$ respectively and points 2 and 4 lie on the same isotherm, the work done by the gas during the cycle is $nR(T_2-T_1)\left(\dfrac{T_2}{T_1}+\dfrac{T_3}{T_2}-\dfrac4a\right)$, then `a' is
\qfig{p1_phy_s3_q17_1}
\qn{18} Following reactions are known as inverse beta decay
\[p+\bar\nu\to n+e^+\qquad n+\nu\to p+e^-\]
These reactions have extremely low probabilities. Because of this, neutrinos and antineutrinos are able to pass through vast amount of matter without any interaction. In an experiment to detect neutrinos, large number of neutrinos coming out from beta decays of a radioactive material were made to pass through a tank of water, containing a cadmium compound in solution, which provided the protons to interact with antineutrinos. Immediately after a proton absorbed a neutrino to yield a positron and a neutron, the positron encountered an electron and both got annihilated. The gamma ray detectors surrounding the tanks responded to the resulting photons. This confirmed that the above reaction has taken place. The maximum energy of this photon is $x$ Mev then $10x$ is (take the mass of an electron to be 0.000537 u.)
\partsec{2}{CHEMISTRY}
\sect{SECTION-I}
\qn{1} Select the correct statements regarding salt solution of weak acid \& strong base
\begin{qoptions}
\opt{A}{Solution will be neutral}
\opt{B}{Degree of hydrolysis increase with increase in dilution}
\opt{C}{Shows anionic hydrolysis}
\opt{D}{Phenolphthalein is suitable indicator during titration of weak acid \& strong base}
\end{qoptions}
\qn{2} Which of the following statements about anhydrous aluminium chloride is/are incorrect:
\begin{qoptions}
\opt{A}{It exists as AlCl$_3$ molecule}
\opt{B}{It is a strong Lewis base}
\opt{C}{It sublims at 100$^\circ$C under vacuum}
\opt{D}{It is not easily hydrolysed}
\end{qoptions}
\qn{3} For the cell\par
Cu $|$ CuSO$_4$ $\|$ HCl $|$ H$_2|$ Pt\par
$E_{cell}$ can be increased by -
\begin{qoptions}
\opt{A}{By adding NH$_3$ in anodic half cell}
\opt{B}{Increasing partial pressure of H$_2$ gas}
\opt{C}{By diluting CuSO$_4$ solution}
\opt{D}{By diluting HCl solution}
\end{qoptions}
\qn{4} Which of the following property of lead is due to the consequence of \textbf{inert pair effect} :
\begin{qoptions}
\opt{A}{Higher stability of Pb$^{4+}$ as compared to Pb$^{2+}$}
\opt{B}{Higher stability of Pb$^{2+}$ as compared to Sn$^{2+}$}
\opt{C}{Size of Pb$^{2+}$ is 133 pm but size of Pb$^{4+}$ is 84 pm}
\opt{D}{Pb$^{4+}$ has oxidising property}
\end{qoptions}
\qn{5} Which of the following are true for an acid-base titration ?
\begin{qoptions}
\opt{A}{Indicators catalyse the acid base reactions by releasing or accepting H$^+$ ions.}
\opt{B}{Indicators do not significantly affect the pH of the solution to which they are added.}
\opt{C}{Acid-base reactions do not occur in absence of indicators.}
\opt{D}{Indicators have different colours in dissociated and undissociated forms.}
\end{qoptions}
\qn{6} Select the correct statement(s)
\begin{qoptions}
\opt{A}{Double chain silicates are known as amphiboles.}
\opt{B}{In cyclic silicates two oxygen atoms per tetrahedron are shared.}
\opt{C}{Orthosilicates contain discrete $\left(\mathrm{SiO_4}\right)^{4-}$ units.}
\opt{D}{Talc and mica are 3D silicates}
\end{qoptions}
\qn{7} $E^\circ_{A^+/A}=0.5$ V ; $E^\circ_{B^+/B}=0.2$ V ; $E^\circ_{C^+/C}=-0.3$ V ; $E^\circ_{D^+/D}=-0.1$ V\par
Based on above data select correct statement under standard condition -
\begin{qoptions}
\opt{A}{Best oxidising agent is $A^+$}
\opt{B}{Best reducing agent is $D$}
\opt{C}{$C$ can displace that $D$ from it's solution}
\opt{D}{In galvanic cell obtained by electrodes $A$ \& $B$ current flow from $B$ to $A$ in external circuit}
\end{qoptions}
\qn{8} Buffer solution $A$ of a weak monoprotic acid and its sodium salt in the concentration ratio $x:y$ has pH $=(\mathrm{pH})_1$. Buffer solution $B$ of the same acid and its sodium salt in the concentration ratio $y:x$ has pH $=(\mathrm{pH})_2$. If $(\mathrm{pH})_2-(\mathrm{pH})_1=1$ unit and $(\mathrm{pH})_1+(\mathrm{pH})_2=9.5$ units, then
\begin{qoptions}
\opt{A}{$\mathrm{p}K_a=4.75$}
\opt{B}{$\dfrac xy=2.36$}
\opt{C}{$\dfrac xy=3.162$}
\opt{D}{$\mathrm{p}K_a=5.25$}
\end{qoptions}
\fwstart
\sect{SECTION-IV}
\qn{9} Match the following\par\nobreak
\begin{center}\begin{tabular}{|L{0.6cm}|L{7.2cm}|L{0.8cm}|L{3.6cm}|}\hline
\multicolumn{2}{|c|}{\textbf{Column I}} & \multicolumn{2}{c|}{\textbf{Column II}}\\\hline
(A) & Negative charge on the anion is equal to the number of terminal oxygen atoms & (P) & $\mathrm{Si_6O_{18}^{-12}}$\\\hline
(B) & Three shared corners per each central atom & (Q) & $\mathrm{SiO_4^{4-}}$\\\hline
(C) & Silicon atom(s) is/are present at the center of geometry and every oxygen atom is present at each corner of the geometry & (R) & $\mathrm{Si_4O_{10}^{-4}}$\\\hline
(D) & Non-planar geometry & (S) & $\mathrm{Si_2O_7^{6-}}$\\\hline
 & & (T) & None of the above mentioned silicate\\\hline
\end{tabular}\end{center}
\qn{10} Match the following\par\nobreak
\begin{center}\begin{tabular}{|L{0.6cm}|L{6.4cm}|L{0.8cm}|L{5cm}|}\hline
\multicolumn{2}{|c|}{\textbf{Column I}} & \multicolumn{2}{c|}{\textbf{Column II} (Suitable indicator at end point)}\\\hline
(A) & Strong acid vs strong base titration & (P) & Methyl orange (3.1 -- 4.4)\\\hline
(B) & Weak acid vs strong base titration & (Q) & Methyl red (4.2 -- 6.3)\\\hline
(C) & Strong acid vs weak base titration & (R) & Phenolphthalein (8.3 -- 10)\\\hline
(D) & Weak acid vs weak base titration & (S) & None of the above mentioned indicator is suitable\\\hline
 & & (T) & All of the above mentioned indicators are suitable\\\hline
\end{tabular}\end{center}
\fwend
\sect{SECTION-II}
\qn{11} In the compound $\mathrm{Na_2\left[B_4O_5(OH)_4\right]\cdot8H_2O}$, if the\par
(i) Number of B--O--B bonds is `$x$'\par
(ii) Number of B--B bonds is `$y$'\par
(iii) Number of sp$^2$ hybridized B atoms is `$z$'\par
Calculate the value of $x+y+z$.
\qn{12} How many litre of 0.3 M NaOH solution should be added to a solution containing 0.5 moles of HCl in hydrogen electrode, to make its Oxidation potential $=0.042$ V at 25$^\circ$C. $\left(\log5=0.7,\ \dfrac{2.303RT}{F}=0.06\text{V}\right)$
\qn{13} 0.1 M of HA is titrated with 0.1 M NaOH, calculate the pH at end point.\par
Given $K_a(\mathrm{HA})=5\times10^{-6}$ and $\alpha\ll1$.
\qn{14} Find number of reagents which produce CO$_2$ when reacted with carbon :\par
dil. H$_2$SO$_4$, conc. H$_2$SO$_4$, dil. HNO$_3$, conc. HNO$_3$, conc. NaOH, Lime water,
\qn{15} Based on the data given at 25$^\circ$C for half cell $\mathrm{M^{(x+n)+},M^{x+}/Pt}$\par
\begin{center}\begin{tabular}{|l|c|c|}\hline
\% Reduced form $(\mathrm{M^{x+}})$ & 50 & 20\\\hline
Potential (V) & 0.1 & 0.112\\\hline
\end{tabular}\end{center}
The value of $n$ is -
\qn{16} The solubility of Pb(OH)$_2$ in water is $5\times10^{-6}$ M. Calculate the pH of buffer solution in which solubility of Pb(OH)$_2$ is $5\times10^{-4}$ mol litre$^{-1}$.
\qn{17} For the cell : Pt$|$Hg($l$)$\,|\,$Hg$_2$Cl$_2$(s)$\,|\,$KCl (1M)$\,\|\,$H$^+|$Q$|$H$_2$Q$|$Pt\par
$E_{cell}=0.1239$ V at 298K, hence pH of cathodic compartment is.\par
Given : $E^\circ_{\mathrm{Calomel,RP}}=0.28$ V, $E^\circ_{\mathrm{Q/H_2Q}}=+0.699$ V
\qn{18} An acidic aqueous solution contains 0.02 M -- FeCl$_2$ and 0.05 M -- FeCl$_3$, at 25$^\circ$C. Now NaOH solution is added is this solution, drop by drop. The pH range in which only the precipitation of Fe(OH)$_3$ occur with out precipitation of Fe(OH)$_2$ differs by `$x$' unit. Find magnitude of `$x$' ?\par
Given : $K_{sp}$ of Fe(OH)$_2=8\times10^{-16}$\par
\hspace*{2.6em}$K_{sp}$ of Fe(OH)$_3=4\times10^{-28}$
\partsec{3}{MATHEMATICS}
\sect{SECTION-I}
\qn{1} If area bounded by the curves $y\le\sqrt{-x},\ x\ge-\sqrt y$ and $y\ge\dfrac14$ is $\dfrac13\left(\dfrac mn\right)^2$, where $m$ \& $n$ are coprime numbers, then which of the following is (are) true-
\begin{qoptions}
\opt{A}{$|m-n|=1$}
\opt{B}{$m$ is prime \& $n$ is even}
\opt{C}{$n$ is prime number and $m$ is even number}
\opt{D}{$m+n$ is an odd number}
\end{qoptions}
\qn{2} The graph of the function $y=f(x)$ passing through the point $(-\pi,1)$ and satisfying the differential equation $\cos^2x\cdot\dfrac{dy}{dx}+y=1+\cos^2x\,e^{-\tan x}$ then-
\begin{qoptions}
\opt{A}{$f\!\left(\dfrac{-\pi}2\right)=1$}
\opt{B}{$f\!\left(\dfrac{\pi^+}2\right)$ approaches to infinity}
\opt{C}{$f'(0)=0$}
\opt{D}{$f'(0)=1-\pi$}
\end{qoptions}
\qn{3} If $\vec r$ is a non-zero vector such that $\vec r\cdot\vec a=\vec r\cdot\vec b=\vec r\cdot\vec c=1$ $(\vec a,\vec b,\vec c$ are noncoplanar) then-
\begin{qoptions}
\opt{A}{If $A(\vec a),B(\vec b),C(\vec c)$ are the vertices of a $\Delta ABC$, then its area is $\left|\dfrac{\vec r}{2}\right|\left[\vec a\ \vec b\ \vec c\right]$}
\opt{B}{If $\vec a,\vec b,\vec c$ are mutually perpendicular unit vectors then $|\vec r|=3$}
\opt{C}{If $A(\vec a),B(\vec b),C(\vec c)$ are vertices of $\Delta ABC$, then its area is $\left|\dfrac{\vec r}{2}\right|$}
\opt{D}{If $\vec a,\vec b,\vec c$ are mutually perpendicular unit vectors then $|\vec r|=\sqrt3$}
\end{qoptions}
\qn{4} If plane $P$ containing lines $L_1:\dfrac{x-1}{k}=\dfrac{y-2}{2}=\dfrac{z-3}{3}$ \& $L_2:\dfrac{x-2}{2}=\dfrac{y-3}{k}=\dfrac{z-4}{3}$, then-
\begin{qoptions}
\opt{A}{$L_1$ \& $L_2$ can be parallel}
\opt{B}{$L_1$ \& $L_2$ can be intersecting}
\opt{C}{Equation of plane $P$ can be $x-y+1=0$}
\opt{D}{Equation of plane $P$ can be $x+y-2z+3=0$}
\end{qoptions}
\qn{5} Let $n$ be a positive integer and $f(n)=1!+2!+3!+\ldots+n!$. $P(x)$ and $Q(x)$ be polynomials in $x$ such that $f(n+2)=P(n)f(n+1)+Q(n)f(n)$ for all $n\ge1$, then
\begin{qoptions}
\opt{A}{$P(x)=x+3$}
\opt{B}{$Q(x)=-x-2$}
\opt{C}{$P(x)=-x-2$}
\opt{D}{$Q(x)=x+3$}
\end{qoptions}
\qn{6} The solution to the differential equation $y\,dx-\left(x+x^3y\right)dy=0$, where $y(1)=-\dfrac32$ can be-
\begin{qoptions}
\opt{A}{$xy^2=\dfrac32$}
\opt{B}{$x^2y=\dfrac32$}
\opt{C}{$xy^2=-\dfrac32$}
\opt{D}{$x^2y=-\dfrac32$}
\end{qoptions}
\qn{7} Let $\vec a,\vec b$ and $\vec c$ be the three vectors having magnitude 1, 1 and 2 respectively. If $\vec a\times\left(\vec a\times\vec c\right)+\vec b=\vec 0$, then (Given that angle between $\vec a$ \& $\vec c$ is acute)
\begin{qoptions}
\opt{A}{$\vec c\cdot\vec b=1$}
\opt{B}{Angle between $\vec a$ and $\vec c$ is $\dfrac\pi6$}
\opt{C}{$\vec a\cdot\vec b=0$}
\opt{D}{$\vec a\times\left(\vec b\times\vec c\right)=\sqrt3\,\vec b$}
\end{qoptions}
\qn{8} Solution of differential equation $\dfrac{x\,dy}{y\,dx}=\dfrac{1+x^2y^2}{1-x^2y^2}$ is-
\begin{qoptions}
\opt{A}{$|y|=k|x|e^{\frac{-x^2y^2}{2}}$}
\opt{B}{$|y|=ke^{\frac{-x^2y^2}{2}}$}
\opt{C}{$|y|=k|x|e^{\frac{x^2y^2}{2}}$}
\opt{D}{$|y|=ke^{\frac{x^2y^2}{2}}$}
\end{qoptions}
\fwstart
\sect{SECTION-IV}
\qn{9} Observe the following column:\par\nobreak
\begin{center}\begin{tabular}{|L{0.6cm}|L{8.4cm}|L{0.8cm}|L{2cm}|}\hline
\multicolumn{2}{|c|}{\textbf{Column I}} & \multicolumn{2}{c|}{\textbf{Column II}}\\\hline
(A) & Area enclosed by the curve $\left(y-\sin^{-1}x\right)^2=x-x^2$ is equal to & (P) & $0$\\\hline
(B) & The area of the region represented by the expression $\sqrt2\le|x+y|+|x-y|\le2\sqrt2$ is equal to & (Q) & $1$\\\hline
(C) & If the area bounded by $y=x^2-3$ and the line $y=ax+2$ attains its minimum value then the parameter `$a$' is equal to & (R) & $\pi$\\\hline
(D) & If $k$ is a positive number and the area of the region bounded by the curves $y=x-kx^2$ and $ky=x^2$ attains its maximum value then $k$ is equal to & (S) & $6$\\\hline
 & & (T) & $\pi/4$\\\hline
\end{tabular}\end{center}
\qn{10} Observe the following columns:\par\nobreak
\begin{center}\begin{tabular}{|L{0.6cm}|L{8.4cm}|L{0.8cm}|L{2cm}|}\hline
\multicolumn{2}{|c|}{\textbf{Column I}} & \multicolumn{2}{c|}{\textbf{Column II}}\\\hline
(A) & If $\vec a+\vec b+\vec c=\alpha\vec d,\ \vec b+\vec c+\vec d=\beta\vec a$ and $\vec a,\vec b,\vec c$ are non-coplanar then the $\left|\vec a+\vec b+\vec c+\vec d\right|$ is & (P) & $\dfrac{2\pi}3$\\\hline
(B) & If $\vec a$ and $\vec b$ are unit vectors inclined at an angle $\theta$ to each other and $\left|\vec a+\vec b\right|<1$, then $\theta$ can be equal to & (Q) & $\dfrac{3\pi}4$\\\hline
(C) & If $\vec a$ is unit vector perpendicular to another unit vector $\vec b$ then $\left|\vec a\times\left[\vec a\times\left\{\vec a\times\left(\vec a\times\vec b\right)\right\}\right]\right|$ is equal to & (R) & $\dfrac{5\pi}6$\\\hline
(D) & Let $\vec a,\vec b,\vec c$ be three unit vectors such that $\vec a+\vec b+\vec c=\vec 0$, then the angle between $\vec a$ and $\vec b$ is equal to & (S) & $0$\\\hline
 & & (T) & $1$\\\hline
\end{tabular}\end{center}
\fwend
\sect{SECTION-II}
\qn{11} Each of the `$n$' urns contains 4 white and 6 black balls. The $(n+1)$th urn contains 5 white and 5 black balls. One of the $(n+1)$ urns is chosen at random and two balls are drawn from it without replacement. Both the balls turn out to be black. If the probability that the $(n+1)$th urn was chosen to draw the balls is $\dfrac1{16}$, then the value of $\dfrac n2$ is equal to
\qn{12} Let a fair coin be tossed 6 times. $A$, $B$ and $C$ three events are defined as\par
$A$ : exactly 4 heads are obtained.\par
$B$ : 4$^{\text{th}}$ head obtained on 5$^{\text{th}}$ toss.\par
$C$ : tail is obtained, on 2$^{\text{nd}}$ or 4$^{\text{th}}$ toss\par
If conditional probability, $P\!\left(\dfrac{B}{A\cap C}\right)=\dfrac mn$ where $m,n\in\mathbb N$ then find the least value of $(n-m)$.
\qn{13} One of the body diagonal of a parallelopiped is $4\hat j-8\hat k$. If two diagonals of two different faces are $6\hat i+6\hat k$ and $-3\hat i+2\hat j-12\hat k$. (given all three diagonals are co-initial). If volume of the parallelopiped be `$v$', then $\left[\sqrt[3]{v}\right]$, (where $[\,.\,]$ represents greatest integer function) is
\qn{14} Let $\vec a=\hat i+\hat j+\hat k,\ \vec b=x_1\hat i+x_2\hat j+x_3\hat k$, where $x_1,x_2,x_3\in\{-3,-2,-1,0,1,2\}$. Number of possible vectors $\vec b$ such that $\vec a$ and $\vec b$ are mutually perpendicular is $\lambda$, then $(\lambda+2)/3$ is
\qn{15} If the planes $x-cy-bz=0,\ cx-y+az=0$ and $bx+ay-z=0$ pass through a straight line, then the value of $a^2+b^2+c^2+2abc$ is:
\qn{16} If three non-zero unequal vectors $\vec\ell=a\hat i+b\hat j+c\hat k,\ \vec m=b\hat i+c\hat j+a\hat k$ \& $\vec n=c\hat i+a\hat j+b\hat k$ are such that $\vec r\cdot\vec\ell=\vec r\cdot\vec m=\vec r\cdot\vec n=0$ where $a,b,c\in R_0$ \& $\vec r$ is non zero vector, then value of $\dfrac{a^2}{bc}+\dfrac{b^2}{ca}+\dfrac{c^2}{ab}$
\qn{17} $P(x,y)$ is a point, which moves in the $xy$ plane such that $2[y]=3[x]$, where $[\,.\,]$ denotes the greatest integer function and $-2\le x\le5$ and $-3\le y\le6$. The area of the region containing the point $P(x,y)$ is equal to
\qn{18} If solution $y=f(x)$ of the differential equation, $x\,dy+y\,dx-y(x+1)\,dx=\left(x^2+x\right)e^{2x}dx$ passes through $(0,0)$, then $\ell nf(1)$ will be
\end{multicols}
\end{document}
